\section*{Bab \ref{ch:g8:midpoints} --- Titik Tengah dan Garis Sejajar}

\begin{solution}{exo:g8:midpoints:1}
$M$ dan $N$ adalah titik tengah dua sisi segitiga $RST$: menurut
teorema titik tengah (\cref{thm:g8:midpoints:direct}),
$(MN) \parallel (ST)$ dan $MN = \frac{ST}{2} = 3.5$ cm.
\end{solution}

\begin{solution}{exo:g8:midpoints:2}
$I$ adalah titik tengah dan $(IJ) \parallel (BC)$: menurut
\emph{konversnya} (\cref{thm:g8:midpoints:converse}), $J$ adalah titik
tengah $[AC]$, jadi $AJ = \frac{11}{2} = 5.5$ cm.
\end{solution}

\begin{solution}{exo:g8:midpoints:3}
Pengukurannya memberi $IJ = 3.5$ cm $= \frac{BC}{2}$. \omterm{def:g6:measure:perimeter}{Keliling}
$AIJ$: $AI = 4$, $AJ = 3$, $IJ = 3.5$: seluruhnya $10.5$ cm --- yaitu
setengah \omterm{def:g6:measure:perimeter}{keliling} $ABC$ ($8 + 6 + 7 = 21$ cm).
\end{solution}

\begin{solution}{exo:g8:midpoints:4}
Tiap sisi $IJK$ adalah setengah sebuah sisi $ABC$: $5$, $6$ dan $8$
cm. \omterm{def:g6:measure:perimeter}{Kelilingnya}: $19$ cm, yaitu setengah $10 + 12 + 16 = 38$ cm.
\end{solution}

\begin{solution}{exo:g8:midpoints:5}
Segitiga titik tengahnya adalah salah satu dari keempat segitiga yang
sama pada \cref{prop:g8:midpoints:medial}: luasnya
$\frac{36}{4} = 9$ cm$^2$ --- dan tiap segitiga kecil dari keempatnya
berluas $9$ cm$^2$.
\end{solution}

\begin{solution}{exo:g8:midpoints:6}
Menurut teorema titik tengah pada $ABC$: $(IJ) \parallel (BC)$, yaitu
$(IJ) \parallel (BK)$, dan $IJ = \frac{BC}{2} = BK$ (karena $K$ titik
tengah $[BC]$). Dua sisi berhadapan $IJKB$ \omterm{def:g6:lines:perp}{sejajar} dan sama panjang:
jadi ia \omterm{def:g7:central:parallelogram}{jajargenjang}
(\cref{thm:g7:central:recognize}, kriteria 3).
\end{solution}

\begin{solution}{exo:g8:midpoints:7}
Jahitannya menghubungkan dua titik tengah: $\frac{6.4}{2} = 3.2$ m.
Pada \omterm{def:g3:division:half}{seperempat} dari \omterm{def:g5:solids:def}{titik sudut} atasnya, gambarnya menyarankan
jahitan sepanjang $\frac{6.4}{4} = 1.6$ m, \omterm{def:g6:lines:perp}{sejajar} dengan alasnya ---
yang dibenarkan Thales pada \cref{ch:g9:thales}.
\end{solution}

\begin{solution}{exo:g8:midpoints:8}
\emph{1.} $M$ adalah titik tengah $[BC]$ dan $(MN) \parallel (AB)$:
menurut konvers teorema titik tengah (pada segitiga $ABC$,
sisi $[AC]$), $N$ adalah titik tengah $[AC]$.

\emph{2.} $(MN) \parallel (AB)$ dan $(AB) \perp (AC)$ (\omterm{def:g3:shapes:rightangle}{sudut siku-siku}
di $A$): garis yang \omterm{def:g6:lines:perp}{sejajar} dengan $(AB)$ juga \omterm{def:g6:lines:perp}{tegak lurus} terhadap
$(AC)$ (\cref{prop:g6:lines:facts}). Jadi $(MN) \perp (AC)$ --- \omterm{def:g6:lines:objects}{ruas
garis} $[MN]$ adalah potongan \omterm{def:g6:symmetry:bisector}{garis sumbu} yang sekali lagi menunjukkan
bahwa $MA = MC$ (\cref{thm:g8:cosine:circle}).
\end{solution}

\begin{solution}{exo:g8:midpoints:9}
\emph{1.} Pada segitiga $ABC$, $I$ dan $J$ adalah titik tengah $[AB]$
dan $[BC]$: $(IJ) \parallel (AC)$ dan $IJ = \frac{AC}{2}$. Pada
segitiga $ACD$, $L$ dan $K$ adalah titik tengah $[DA]$ dan $[CD]$:
$(LK) \parallel (AC)$ dan $LK = \frac{AC}{2}$.

\emph{2.} Jadi $[IJ]$ dan $[LK]$ \omterm{def:g6:lines:perp}{sejajar} (keduanya \omterm{def:g6:lines:perp}{sejajar} dengan
diagonal $(AC)$) dan sama panjang (keduanya setengah $AC$): $IJKL$
adalah \omterm{def:g7:central:parallelogram}{jajargenjang} (\cref{thm:g7:central:recognize}, kriteria 3) ---
untuk \emph{setiap} \omterm{def:g4:geometry:polygon}{segi empat} $ABCD$, seperti yang diduga pada
\cref{exo:g7:central:11}.
\end{solution}

\begin{solution}{pb:g8:midpoints:1}
\textbf{1.} Segitiga apa pun yang digambar, ketiga \omterm{pb:g8:midpoints:1}{garis beratnya}
tampak melewati satu titik, yang terletak di dalam segitiganya, jelas
lebih dekat ke titik tengah tiap sisinya daripada ke \omterm{def:g5:solids:def}{titik sudut} yang
di seberangnya. Sisa soalnya membuktikannya.

\textbf{2.} Ambillah \omterm{pb:g8:midpoints:1}{garis berat} $[AK]$ pada segitiga $ABC$: segitiga
$ABK$ dan $AKC$ punya alas $BK$ dan $KC$ yang sama panjang
($K$ adalah titik tengah $[BC]$) yang dibawa garis yang sama
$(BC)$, dan tinggi yang sama: yaitu jarak dari $A$ ke $(BC)$.
Menurut rumus luasnya, masing-masing berluas
alas $\times$ tinggi $\div\ 2$, dengan alas yang sama dan tinggi yang
sama: jadi luasnya sama. Untuk segitiga seluas
$36$ cm$^2$: dua keping seluas $18$ cm$^2$ masing-masing.

\textbf{3.} $I$ dan $J$ adalah titik tengah sisi $[AB]$ dan
$[AC]$ pada segitiga $ABC$, jadi menurut
\cref{thm:g8:midpoints:direct}: $(IJ) \parallel (BC)$ dan
$IJ = \frac{BC}{2}$.

\textbf{4.} $M$ dan $N$ adalah titik tengah sisi $[GB]$ dan
$[GC]$ pada segitiga $GBC$, jadi menurut teorema yang sama:
$(MN) \parallel (BC)$ dan $MN = \frac{BC}{2}$.

\textbf{5.} Baik $(IJ)$ maupun $(MN)$ \omterm{def:g6:lines:perp}{sejajar} dengan $(BC)$, jadi
keduanya \omterm{def:g6:lines:perp}{sejajar} satu sama lain; dan $IJ = \frac{BC}{2} = MN$. Karena
itu \omterm{def:g4:geometry:polygon}{segi empat} $IJNM$ punya dua sisi berhadapan, $[IJ]$ dan
$[NM]$, yang \omterm{def:g6:lines:perp}{sejajar} dan sama panjang: menurut kriteria 3 pada
\cref{thm:g7:central:recognize}, $IJNM$ adalah \omterm{def:g7:central:parallelogram}{jajargenjang}.

\textbf{6.} $G$ terletak pada \omterm{pb:g8:midpoints:1}{garis berat} $[CI]$, dan $N$ adalah
titik tengah $[GC]$, yaitu \omterm{def:g6:lines:objects}{ruas garis} yang dibawa \omterm{pb:g8:midpoints:1}{garis berat} yang
sama: jadi $I$, $N$ (dan $G$, $C$) semuanya terletak pada $(CI)$.
Begitu pula $J$, $M$ dan $G$ terletak pada \omterm{pb:g8:midpoints:1}{garis berat} $(BJ)$. Pada
\omterm{def:g7:central:parallelogram}{jajargenjang} $IJNM$, diagonalnya adalah $[IN]$ dan $[JM]$; keduanya
dibawa kedua \omterm{pb:g8:midpoints:1}{garis berat} itu, yang berpotongan di $G$ --- jadi
diagonalnya berpotongan di $G$. Menurut \emph{definisi} \omterm{def:g7:central:parallelogram}{jajargenjang}
(kriteria 1 pada \cref{thm:g7:central:recognize}: diagonalnya punya
titik tengah yang sama), $G$ adalah titik tengah $[IN]$ \emph{dan}
$[JM]$: jadi $GI = GN$ dan $GJ = GM$.

\textbf{7.} $M$ adalah titik tengah $[GB]$, jadi $BM = MG$; dan
$MG = GJ$ menurut pertanyaan 6. Karena itu $BM = MG = GJ$: potongan
\omterm{pb:g8:midpoints:1}{garis berat} $[BJ]$ terpotong menjadi tiga bagian yang sama, yang $BG$ mengambil dua:
\[
BG = 2 \times GJ .
\]
Secara \omterm{def:g7:central:def}{simetris}, $N$ adalah titik tengah $[GC]$ dan $GN = GI$, jadi
$CN = NG = GI$ dan $CG = 2 \times GI$: pada \omterm{pb:g8:midpoints:1}{garis berat} $[CI]$ pun,
titik potongnya duduk pada dua pertiga jalan dari \omterm{def:g5:solids:def}{titik sudutnya}.

\textbf{8.} Alasannya tidak memakai apa pun yang khusus tentang
pasangan $[BJ]$, $[CI]$: kalau dijalankan pada pasangan $[BJ]$,
$[AK]$, alasan itu menunjukkan bahwa titik potongnya juga terletak
pada $[BJ]$ pada dua pertiga jalan dari
$B$. Tetapi hanya ada \emph{satu} titik pada \omterm{def:g6:lines:objects}{ruas garis} $[BJ]$ yang
berjarak $\frac23 BJ$ dari $B$ --- jadi titik potong
$[BJ]$ dan $[AK]$ adalah $G$ yang persis sama. Karena itu ketiga
\omterm{pb:g8:midpoints:1}{garis beratnya} melewati $G$: ketiganya \emph{konkuren}.

\textbf{9.} $BG = \frac23 \times 9 = 6$ cm dan
$GJ = \frac13 \times 9 = 3$ cm.

\textbf{10.} Menurut pertanyaan 8, perhitungan dua pertiga yang sama
berlaku pada \omterm{pb:g8:midpoints:1}{garis berat} $[AK]$ (jalankan alasan pertanyaan 4--7
dengan pasangan $[AK]$, $[BJ]$): $AG = 2 \times GK$.

\textbf{11.} Pada segitiga $GBC$, titik $K$ adalah titik tengah
sisi $[BC]$, jadi $[GK]$ adalah \omterm{pb:g8:midpoints:1}{garis berat} $GBC$; menurut pertanyaan
2, segitiga $GBK$ dan $GKC$ punya luas yang sama.

\textbf{12.} \omterm{pb:g8:midpoints:1}{Garis berat} $[AK]$ memotong $ABC$ menjadi $ABK$ dan
$AKC$ yang luasnya sama (pertanyaan 2). Dengan membuang dari
masing-masingnya segitiga kecil pada pertanyaan 11:
\[
\text{luas}(ABG) = \text{luas}(ABK) - \text{luas}(GBK),
\qquad
\text{luas}(ACG) = \text{luas}(AKC) - \text{luas}(GKC),
\]
dan luas yang dibuang itu sama, jadi
$\text{luas}(ABG) = \text{luas}(ACG)$. Perhitungan yang sama
sepanjang \omterm{pb:g8:midpoints:1}{garis berat} $[BJ]$ --- yang membagi dua $ABC$ menjadi $ABJ$
dan $CBJ$, sedangkan $[GJ]$, \omterm{pb:g8:midpoints:1}{garis berat} segitiga $GAC$, membaginya
menjadi $GAJ$ dan $GCJ$ --- memberi
$\text{luas}(ABG) = \text{luas}(CBG)$. Jadi ketiga segitiga
$ABG$, $BCG$, $CAG$ punya luas yang sama, dan karena bersama-sama
ketiganya mengubin $ABC$, masing-masing berluas sepertiga seluruhnya.

\textbf{13.} Pada segitiga $ABG$, titik $I$ adalah titik tengah
sisi $[AB]$, jadi $[GI]$ adalah \omterm{pb:g8:midpoints:1}{garis berat} $ABG$ \emph{dari \omterm{def:g5:solids:def}{titik
sudut} $G$}: ia memotong $ABG$ menjadi dua segitiga yang luasnya sama
(pertanyaan 2), yaitu $AIG$ dan $IBG$. Hal yang sama terjadi pada
$BCG$ (\omterm{pb:g8:midpoints:1}{garis berat} $[GK]$) dan pada $CAG$ (\omterm{pb:g8:midpoints:1}{garis berat} $[GJ]$). Tiap
sepertiga $ABC$ terpotong menjadi dua: keenam irisan di sekeliling
$G$ masing-masing berluas
$\frac13 \times \frac12 = \frac16$ segitiganya.

\textbf{14.} Pada \omterm{pb:g8:midpoints:1}{garis berat} $[AK]$, dari $A(0,0)$ ke $K(3,3)$: $G$
duduk pada dua pertiga jalannya, jadi koordinatnya adalah
$\left(\frac23 \times 3,\ \frac23 \times 3\right) = (2, 2)$. Pada
\omterm{pb:g8:midpoints:1}{garis berat} $[BJ]$, dari $B(6, 0)$ ke $J(0, 3)$: dua pertiga
jalannya berarti menambahkan dua pertiga perpindahannya ($-6$
mendatar, $+3$ tegak):
\[
\left(6 + \tfrac23 \times (-6),\ 0 + \tfrac23 \times 3\right)
= (2, 2) .
\]
Titik yang sama --- seperti yang dijanjikan pertanyaan 8. Dan
$\frac{0 + 6 + 0}{3} = 2$, $\frac{0 + 0 + 6}{3} = 2$: koordinat \omterm{pb:g8:midpoints:1}{titik
beratnya} adalah \emph{rata-rata} koordinat ketiga
\omterm{def:g5:solids:def}{titik sudutnya}.

\textbf{15.} Menurut pertanyaan 2, mata pisau yang dibaringkan
sepanjang sebuah \omterm{pb:g8:midpoints:1}{garis berat} punya karton yang sama banyak pada tiap
sisinya (luasnya sama), jadi segitiganya seimbang sepanjang
masing-masing dari ketiga \omterm{pb:g8:midpoints:1}{garis beratnya}. Titik keseimbangannya harus
terletak pada ketiga garis keseimbangannya sekaligus, dan menurut
pertanyaan 8 ketiga \omterm{pb:g8:midpoints:1}{garis berat} itu berbagi tepat satu titik:
yaitu $G$. Di situlah ujung pensilnya diletakkan. (Luas yang sama pada
kedua sisinya adalah versi alasan yang masuk akal; versi yang jujur,
yang menimbang \emph{di mana} kartonnya duduk dan bukan hanya berapa
banyaknya, perlu hitung integral pada volume di tingkat universitas
--- yang membenarkan $G$.)

\end{solution}
